\clearpage
\fancyhead[L]{\small BODE. ASTRONOMISCHES JAHRBUCH.}
\hypertarget{native-624-lower}{}
\begin{center}\small Göttingische gelehrte Anzeigen. Stück 23. Seite 225..232. 1820 Februar 7.\\[3pt]\rule{.72\linewidth}{.3pt}\end{center}

\emph{Astronomisches Jahrbuch für das Jahr {\upshape 1822}, nebst einer Sammlung der neuesten in die astronomischen Wissenschaften einschlagenden Abhandlungen, Beobachtungen und Nachrichten. Berechnet und herausgegeben von} Dr. J. E. \textsc{Bode}, \emph{Kön. Astronomen u.s.w.} Berlin. Bei dem Verf. und in Comm. bei F. Dümmler. 260 Seiten in 8. nebst einer Kupfertafel.

\begingroup\rightskip=0pt plus 1em
Von zwei Sonnen- und zwei Mondfinsternissen, die sich im Jahr 1822 ereignen, ist in Europa nur die eine Mondfinsterniss sichtbar; auch der am 5. November statt findende Durchgang des Mercur kann in Europa nicht beobachtet werden. Die Einrichtung des Jahrbuchs ist dieselbe geblieben, wie bei den vorhergehenden Jahrgängen; nur die Jupiterstrabanten-Verfinsterungen sind, zur Erleichterung der Arbeit, nach den \textsc{Wargentin}'schen Tafeln berechnet, und blos bei den ersten Trabanten hin und wieder die \textsc{Delambre}'schen Tafeln mit zu Hülfe genommen.\par\endgroup

Unter den Abhandlungen, welche der Anhang enthält, macht den Anfang ein Versuch, die astrognostischen Namen einzelner Sterne zu befestigen, von Hrn. Prof. \textsc{Buttmann}. Dieser gelehrte Philologe \SourceLeaf{625} stellt hier ein Verzeichniss von 80 Sternnamen auf, grösstentheils aus der Arabischen Sprache entlehnt, von denen einige schon sonst gangbar oder doch nicht unbekannt, andere.berichtigt oder an eine andere Stelle gesetzt, und noch andere ganz neu gebildet sind (z. B. Clavus für $\beta$ im Stier, Meion für $\varkappa$ im Orion, 
Sectatrix für $\beta$ im Fuhrmann). Bei den Astronomen ist eigentlich nur eine sehr kleine Anzahl von eignen Namen für die Fixsterne \emph{wirklich gangbar} (etwa ein Dutzend), und die Versuche einiger, ältere Arabische Namen wieder einzuführen, haben keinen Beifall gefunden. Wollte man aber doch einmal eine grössere Anzahl in Gebrauch bringen, so würde es wohl am zweckmässigsten sein, sich schlechtweg und \emph{ohne alle Abänderung} an die Namen zu halten, die \textsc{Piazzi} in seinen neuen Sterncatalog aufgenommen hat. Jede willkürliche Abänderung könnte, statt zur Verbesserung, nur zur Sprachverwirrung führen. --- Jupiterstrabanten-Verfinsterungen, Sternbedeckungen, Mondfinsternisse und Oppositionen der Vesta, des Uranus, Jupiter, Saturn und der Pallas im Jahre 1818 beobachtet von Hrn. Prof. \textsc{Bürg}. --- Algols Lichtperiode aus neuern Beobachtungen bestimmt und Berechnung des kleinsten Lichts desselben für 1820 bis 1822 vom Hrn. Prof. \textsc{Wurm}. Diese Untersuchung zeigt, dass die Länge der Periode, seit der Entdeckung der Veränderlichkeit des Sterns, durchaus keine merkliche Veränderung erlitten hat. --- Astronomische Beobachtungen auf der Prager Sternwarte, von den Herren \textsc{David} und \textsc{Bittner}. --- Noch von Hrn. Prof. \textsc{Bürg} Beobachtungen der Juno, und Sternbedeckungen 1819, nebst astronomischen Nachrichten. Die hier vorkommende Verantwortung gegen einen Angriff oder eine falsche Auslegung eines Zeitungsartikels, den einen Cometen von 1818 betreffend, welche sich in der in Genua herauskommenden Correspondance astronomique befindet, veranlasst den Wunsch, dass Persönlichkeiten von astronomischen Verhandlungen immer entfernt bleiben mögen. --- Beobachtete Sternbedeckungen und Jupiterstrabanten-Verfinsterungen in Prag von Hrn. Prof. \textsc{Hallaschka}. --- Aeltere Beobachtungen des Uranus von \textsc{Lemonnier} (s. Jahrg. 1819 S. 1909). --- Ueber die Sonnenfinsterniss vom 7ten Sept. 1820 für Deutschland und die angrenzenden Länder von Hrn. Prof. \textsc{Littrow}. Es wird hier Anfang und Ende für 230 Oerter in Deutschland und den benachbarten Ländern, auch die nördliche und südliche Grenzlinie der Zone, innerhalb welcher die Finsterniss ringförmig erscheinen wird, durch eine Anzahl Punkte bestimmt, angegeben. Der grössere Theil von Deutschland liegt innerhalb dieser Zone; die südliche Grenze geht nahe bei Aachen, Strassburg, Freiburg und Zürich, die nördliche nahe bei Schwerin, Wittenberg und Tabor in Böhmen vorbei. Es wird sehr wünschenswerth sein, dass die Dauer der ringförmigen Finsterniss an recht vielen Orten, besonders solchen, die den Grenzen nahe liegen, beobachtet werde, wozu nichts weiter als ein gutes Fernrohr und eine Secundenuhr erforderlich ist. --- Beobachtungen der Vesta, Pallas, des Uranus, Saturn und Jupiter, und der Sonnenfinsterniss im Jahr 1818, zu Wilna von Hrn. Prof. \textsc{Sniadecky}. --- Ueber die geographische Länge der Berliner Sternwarte. Diese Länge ist früher immer zu den sehr schwankenden gezählt; allein seit dem die Sternwarte mit bessern Hülfsmitteln zur Zeitbestimmung versehen ist, geben auch die Beobachtungen eine mehr befriedigende Uebereinstimmung. Aus 23 Sternbedeckungen und 4 Sonnenfinsternissen, die von 1802 bis 1815 beobachtet sind, .folgt im Mittel die Länge $44^{\mathrm{m}}10^{\mathrm{s}}5$ von Paris. --- Beobachtete Gegenscheine des Mars, Uranus und Jupiter 1817, 1818, ferner der Sonnenfinsterniss von 1818 und des \textsc{Enke}'schen Cometen, von Hrn. \textsc{Derflinger} in Kremsmünster. --- Beobachtung einer Bedeckung des Mars vom Monde, in Wien von Hrn. Prof. \textsc{Bürg}. --- Astronomische Beobachtungen auf der Berliner Sternwarte im Jahr 1818. Eine Berichtigung bedarf die hier vorkommende übereilte Behauptung, dass die Bestimmung des heliocentrischen Orts aus einem nahe bei der Opposition beobachteten geocentrischen wegen der Kleinheit der Winkel des Dreiecks unsicher werde, und dass es sogar vortheilhafter sei, von der Opposition entferntere Beobachtungen zum Grunde zu legen. --- Beobachtungen eines Cometen 1818 und 1819 von Hrn. Prof. \textsc{Bessel}; ferner Planetenoppositionen, Sonnenfinsterniss und Sonnenwenden. --- Von Hrn. Dr. \textsc{Olbers} über die höchst merkwürdige \SourceLeaf{626} von Hrn. Prof. \textsc{Enke} gemachte, und in unsern Blättern (St. 28 u. 83 v. v. J. [S. 417 u.f. d. B.]) zuerst angezeigte Entdeckung der Identität des einen Cometen von 1805 mit dem einen von 1818. Hr. Dr. \textsc{Olbers} fügte dieser grossen Entdeckung die nicht weniger wichtige Bemerkung bei, dass theils auch der Comet von 1795 höchst wahrscheinlich derselbe Comet gewesen ist, theils auch die in der Connaissance des tems 1819 bekannt gemachten zwei Cometen-Beobachtungen von 1786 eben demselben Cometen angehören. Ferner theilt hier dieser treffliche Astronom seine ersten Beobachtungen des im Juli d. J. erschienenen Cometen mit. --- Der folgende schöne Aufsatz des Hrn. Prof. \textsc{Enke} gibt nun einen ausführlichen Bericht über seine Arbeit, jenen Cometen von kurzer Umlaufszeit betreffend, und erhebt die obige wahrscheinliche Vermuthung des Hrn. Dr. \textsc{Olbers} zur moralischen Gewissheit. In der Mitte Mais 1822 wird nun dieser Comet abermals in seine Sonnennähe kommen; allein unglücklicherweise wird er vorher, seiner Stellung gegen die Erde wegen, nur äusserst lichtschwach sein, so dass es Noth haben wird, ihn bis dahin zu beobachten. Nach dem Durchgang durch die Sonnennähe hingegen wird er zwar hell genug sein, allein zu weit südlich stehen, um in Europa sichtbar zu werden. Mit Leichtigkeit aber wird er alsdann auf der südlichen Hälfte der Erdkugel beobachtet werden können, wenn es nur nicht dort an Beobachtern fehlen wird. Bei Gelegenheiten, wie diese, eben so wie bei so vielen andern, macht sich das Bedürfniss einer \emph{südlichen} Sternwarte recht fühlbar, und wir fühlen uns bei dieser Veranlassung gedrungen, unsre innige Ueberzeugung auszusprechen, dass nach den grossen Unterstützungen, welche die Astronomie in den neuesten Zeiten in Europa erhalten hat, dieser Wissenschaft kein grösserer Dienst mehr geleistet werden könnte, als durch Errichtung einer wohlausgerüsteten und mit einem einsichtsvollen, erfahrnen und thätigen Astronomen zu besetzenden Sternwarte in der südlichen Hemisphäre, etwa auf dem Vorgebirge der guten Hoffnung. --- Von demselben Astronomen Beobachtungen des Cometen von 1819, und Elemente seiner Bahn, ferner Beobachtungen des Uranus, Saturn, Jupiter, der Ceres und Juno 1818 und 1819. --- Auflösung der Aufgabe, die Declination eines Gestirns ohne Winkelinstrument blos vermittelst eines Fernrohrs zu finden, von Hrn. Prof. \textsc{Fischer} (aus der Dauer eines solchen Durchganges durch das Gesichtsfeld, wobei Eintritt und Austritt an den Endpunkten eines Durchmessers geschehen). --- Geometrischer Lauf der Ceres 1819 und 1820, berechnet von Hrn. Dr. \textsc{Westphal}; und geocentrischer Lauf der Juno 1820 berechnet von Hrn. Prof. \textsc{Nicolai} --- Beobachtungen des Cometen vom Juli 1819 auf der Berliner Sternwarte (gehen vom 2. bis 27. Juli). --- Von Hrn. Prof. \textsc{Nicolai}, Beobachtungen der Juno auf der Mannheimer Sternwarte, und neue Elemente der Bahn aus diesen und Göttinger Beobachtungen abgeleitet. Ferner die ersten Beobachtungen mit dem wiederhergestellten dreifussigen \textsc{Reichenbach}'schen Multiplicationskreise vom Nordstern, welche mit einer sehr schönen Uebereinstimmung die Polhöhe der Mannheimer Sternwarte $49^\circ\,29'\,12''95$ geben. Endlich Beobachtungen und Elemente des Cometen von 1819. --- Die hierauf folgenden Bemerkungen über den Gebrauch der Libelle und des Loths zur Rectification astronomischer Werkzeuge von einem Ungenannten enthalten nichts neues, und beweisen nur, dass es nicht gleichgültig ist, auf welche Art Libellen an einem Instrument angebracht sind, dass sie einen ihrer Bestimmung angemessenen Grad von Empfindlichkeit haben müssen, und dass sie einen vorsichtigen Beobachter erfordern. Der hier gemachte Vorschlag, Veränderlichkeit oder Unveränderlichkeit der Richtung der Schwere durch ein in einem Brunnen aufgehängtes Loth zu prüfen, möchte schwerlich zu etwas anderm führen, als die \emph{Unerreichbarkeit absoluter Festigkeit} bei jeder Aufstellung von Messungsapparaten zu beweisen. --- Auszug aus einem Schreiben des Hrn. General von \textsc{Lindener} aus Glatz. Da bei dem im Juli d. J. erschienenen Cometen durch einen glücklichen Zufall die untere Conjunction Vormittags am 26. Juni mit dem Durchgang durch den aufsteigenden Knoten zusammenfiel, so war es sehr erwünscht, dass dieser eifrige Liebhaber der Astronomie gerade zu dieser Zeit die Sonne wirklich beobachtet hatte, und so das durch wirkliche Erfahrung \SourceLeaf{627} bestätigen konnte, was man auch, nach unsern übrigen gegenwärtigen Kenntnissen, nicht anders erwarten konnte, dass nemlich von dem lockern Cometenkörper, obgleich er gewiss und genau zwischen der Erde und Sonne stand, doch auf der Sonnenscheibe auch nicht die allergeringste Spur erkannt wurde. Wir erinnern uns in öffentlichen Blättern gelesen zu haben, dass noch ein anderer Beobachter in Wien, welcher in derselben Stunde gleichfalls die Sonne beobachtete, gar nichts besonderes wahrgenommen hat. --- Beobachtungen dieses Cometen in der letzten Hälfte des August, von Hrn. Dr. \textsc{Olbers}, wie auch Elemente seiner Bahn von Hrn. \textsc{Bouvard} und von Hrn. \textsc{Dirksen}. --- Ueber die Aufgabe, den Ort eines Gestirns aus beobachteten Allignements zu finden, von Hrn. Dr. \textsc{Olbers}. Im Jahrbuche für 1821 findet sich eine Auflösung dieser Aufgabe von \textsc{Bessel}: bei der von \textsc{Olbers} gewählten Einkleidung braucht man noch zwei Logarithmen weniger. Doch muss erinnert werden, dass die zu dieser erforderlichen 20 Logarithmen an eben so viel verschiedenen Stellen stehen, während die bei der \textsc{Bessel}'schen Form nöthigen 22 nur an 17 verschiedenen Stellen aufzuschlagen sind. Auch bemerken wir noch, dass, wenn man die bekannte Hülfstafel für Logarithmen von Summen und Differenzen mit gebraucht, man mit 18 Logarithmen an 16 Stellen, oder mit 20 Logarithmen an 14 Stellen ausreichen kann. --- Beobachtungen des Cometen vom Juli 1819, Beschreibung des sechsfussigen \textsc{Reichenbach}'schen Mittagsfernrohrs; Beobachtung der Jupiters-Opposition, und des Polarsterns, von Hrn. Hofr. \textsc{Gauss}. Wir machen, bei dem, was hier über die absolute Zuverlässigkeit der Zenithdistanzen gesagt ist, auf zwei den Sinn ganz verfälschende Druckfehler aufmerksam; S. 239 [S. 427] oben muss nemlich statt merklich, \emph{unmerklich}, und in der 5ten Zeile von unten muss statt doch kann ich schon glauben, gelesen werden, \emph{doch kann ich schwer glauben}. --- Beobachtungen des Cometen vom Juli 1819 auf der Krakauer Sternwarte, von Hrn. Prof. \textsc{Leski}. --- Verzeichniss der genauen Länge und Breite von neun der vornehmsten Fixsterne, von denen die Mondsabstände in dem Nautical Almanac und in der Connoissance des tems vorkommen, aus letzterer für 1821 entlehnt (s. vorigjähr. Anz. S. 1907 [S. 615 d. B.]). --- Marseiller Beobachtungen eines neuen Cometen, und dessen Elemente der Bahn berechnet von Hrn. Prof. \textsc{Enke}. Dieser Comet wurde, so viel bisher bekannt ist, blos in Marseille im Juni 1819 beobachtet, und es ist sehr merkwürdig, dass sich die Beobachtungen weit besser in einer Ellipse von einer $2\frac{1}{2}$jährigen 
Umlaufszeit, als in einer Parabel darstellen lassen. --- Beobachtungen des (andern) Cometen von 1819 auf der Greenwicher Sternwarte, und daraus von Hrn. \textsc{Rümker} abgeleitete parabolische Elemente. --- Die zum Schluss noch folgenden kurzen Nachrichten enthalten noch mehrere astronomische Beobachtungen und Notizen, die hier nicht besonders angeführt werden können. Auf der beigefügten Kupfertafel ist der geocentrische Lauf des Cometen vom Juli 1819, und eine Projection seiner wahren Bahn, wie auch die des \textsc{Enke}'schen und des andern vom November 1818 bis Januar 1819 beobachteten Cometen vorgestellt.
